\scanpage{98}
\mspage{6.}
% archive sheet number: 97
2°) \fbox{$X_0-D'_0=\mathbb P^1_k$} i.e. $X_0=\mathbb P^1_k$, $D'_0=\varnothing$. Alors
$\deleted{\pi_1^t(\unsure{X|\mathfrak X}/D,\bar\xi_1)}\inserted{\pi_1^t(\mathfrak X/D,\bar\xi_1)}$ est \deleted{une extension} isomorphe à
$\mu_n(k)$, \deleted{\unsure{où}} $n$ est la partie première à
$p=\mathrm{car}\,k$ du degré de la self-intersection $X_0.X_0$ dans $\mathfrak X$.

\uline{Exemple} II a) $X_0=\operatorname{Spec}\mathbb Z$. \inserted{[plus gén. dans $X_0$ \unsure{ttes|tous} les \unsure{car.} résiduelles existent]} Alors $\mu^\infty_{t,X_0}=0$,
d'où $\pi^t_2(\mathfrak X/D,\bar\xi_1)\simeq\pi_1^t(X_0/D'_0,\bar\xi_0)$.

\deleted{\unsure{Le Pb de}} Si $D'_0=\varnothing$, on trouve $0$.

Si $D'_0\ne\varnothing$, la partie abélienne de
$\pi_1^t(X_0/D'_0,\bar\xi_0)$ est calculable par le corps de
classes \inserted{$(\simeq\prod_{\ell\in D'_0}(\mathbb Z/\ell\mathbb Z)^*)$} ; sans doute $\pi_1^t(X_0/D'_0,\bar\xi_0)$ n'est
en général pas abélien \unsure{[i.e.} il y a des \margin{?}
\deleted{\unsure{revêtements}} corps de \unsure{nbres|nombres} tamely ramifiés
sur $\unsure{Q|\mathbb Q}$, qui ne sont pas abéliens \ldots. ??].

b) $X_0=\operatorname{Spec}\mathbb Z-S$, $S\ne\varnothing$. Alors
$\mu^\infty_{t,X_0}=\prod_{\ell\in S}\mu^\infty(\ell)_{X_0}$. Mais \unsure{qui est|qu'est} $\pi^t_2(X_0/D'_0)$

???
